\documentclass[11pt]{article}
\usepackage[T1]{fontenc}
\usepackage{lmodern}
\usepackage{amsmath,amssymb,amsthm}
\usepackage[margin=1in]{geometry}
\usepackage[hidelinks]{hyperref}

\newtheorem{theorem}{Theorem}
\newtheorem{lemma}[theorem]{Lemma}
\newtheorem{corollary}[theorem]{Corollary}
\newcommand{\Z}{\mathbb Z}
\newcommand{\inj}{\operatorname{Inj}}
\newcommand{\Tot}{\operatorname{Tot}}

\title{Torsion-free integral homology for anchored configuration spaces of finite simple graphs}
\author{Xiaojian Zeng\\
PBC School of Finance, Tsinghua University, 100190\\
\texttt{xiaojianzeng@foxmail.com}}
\date{}

\begin{document}
\maketitle

\begin{abstract}
For a finite graph $G$, a vertex set $K\subseteq V(G)$, and $n\geq |K|$, the anchored configuration space $\Sigma(G,K,n)$ is the subcomplex of $G^n$ in which every anchor occurs as a coordinate.  We prove that $H_i(\Sigma(G,K,n);\Z)$ is free abelian for every $i$.  For connected $G$ we obtain an additive decomposition in terms of tensor powers of $H_1(G;\Z)$ and the homology of a compatibility nerve.  A cellular Cech resolution identifies that nerve with the corresponding anchored space of a tree, so its homology is free by Kozlov's wedge-of-spheres theorem.  The disconnected case follows by a component decomposition and the integral K\"unneth theorem.  In particular, anchored configuration spaces of cographs have no integral homology torsion, resolving the associated existence question in the negative.
\end{abstract}

\section{Statement and context}

Regard a finite simple graph $G$ as a one-dimensional CW complex, with the product cubical structure on $G^n$.  For $K\subseteq V(G)$ and $n\geq |K|$, set
\[
 \Sigma(G,K,n)=\{(x_1,\ldots,x_n)\in G^n:K\subseteq\{x_1,\ldots,x_n\}\}.
\]
Thus the coordinates are labelled and collisions are allowed.  This is the anchored configuration space of Kozlov and of Definition~4.13 in Mamun--Nalikka--Ramos~\cite{kozlov2023,mamun2026}.  Kozlov determined the homotopy type for trees and later computed the homology of generalized anchored spaces on circle graphs~\cite{kozlov2023,kozlov2024}.  Mamun--Nalikka--Ramos proved finite-generation and uniform-exponent statements for anchored spaces of cographs and asked whether any anchored configuration space has integral homology torsion~\cite[Corollary~4.18 and p.~28]{mamun2026}.

We answer that existence question negatively, in the stronger class of all finite graphs.  No priority claim beyond the cited and bounded comparison is intended.

\begin{theorem}\label{thm:main}
Let $G$ be a finite simple graph, let $K\subseteq V(G)$, and let $n\geq |K|$.  Then
\[
 H_i(\Sigma(G,K,n);\Z)
\]
is free abelian for every $i\geq0$.
\end{theorem}

The proof gives more information when $G$ is connected.  Put $r=|K|$ and $L=H_1(G;\Z)$.  For integers $m\geq r\geq1$, define a simplicial complex $N_{r,m}$ as follows.  Its vertices are the injections from a fixed $r$-element set $K$ to $[m]=\{1,\ldots,m\}$.  A set $\{f_0,\ldots,f_p\}$ is a simplex precisely when its requirements are compatible:
\begin{equation}\label{eq:compatibility}
 f_a(k)=f_b(k')\quad\Longrightarrow\quad k=k'
 \qquad(0\leq a,b\leq p;\ k,k'\in K).
\end{equation}
The isomorphism type depends only on $(r,m)$.

\begin{theorem}[Additive decomposition]\label{thm:decomposition}
If $G$ is connected and $r=|K|\geq1$, then, for every $i$,
\begin{equation}\label{eq:decomposition}
 H_i(\Sigma(G,K,n);\Z)\cong
 \bigoplus_{\substack{J\subseteq[n]\\ |J|\leq n-r}}
 L^{\otimes |J|}\otimes
 H_{i-|J|}(N_{r,n-|J|};\Z),
\end{equation}
where homology in negative degrees is zero.  This is an isomorphism of abelian groups; no naturality with respect to the choice of a basis of $L$ is asserted.
\end{theorem}

\section{The anchor-injection resolution}

Assume throughout this section that $G$ is connected and $K\neq\varnothing$.  For every $f\in\inj(K,[n])$, let
\[
 U_f=\{(x_1,\ldots,x_n)\in G^n:x_{f(k)}=k\text{ for all }k\in K\}.
\]
These cubical subcomplexes cover $\Sigma(G,K,n)$.  A finite intersection is empty exactly when two of its requirements assign distinct anchors to the same coordinate.  Otherwise, if $S\subseteq[n]$ is the set of constrained coordinates, the intersection is canonically $G^{[n]\setminus S}$.

Fix an order on $\inj(K,[n])$ and form the cellular Cech double complex
\begin{equation}\label{eq:cech}
 D_{p,q}=\bigoplus_{f_0<\cdots<f_p}
 C_q(U_{f_0}\cap\cdots\cap U_{f_p};\Z),
\end{equation}
where empty intersections are omitted.  The horizontal differential is the alternating sum of inclusions, and the vertical differential is the cellular boundary, with the usual total-complex sign convention.

\begin{lemma}\label{lem:augmentation}
The augmentation $\Tot D\to C_*(\Sigma(G,K,n);\Z)$ is a quasi-isomorphism.
\end{lemma}

\begin{proof}
This is a cellwise calculation, so no open-cover hypothesis is involved.  Fix a product cell $c$ of the union.  The horizontal summands having basis element $c$ form the augmented simplicial chain complex on the nonempty set
\[
 \{f\in\inj(K,[n]):c\subseteq U_f\}.
\]
Every subset of this index set has nonempty intersection because all its requirements already hold on $c$.  The resulting augmented simplex is exact.  Decomposing the horizontal augmentation by product cells proves the assertion.
\end{proof}

Choose a spanning tree $T\subseteq G$, orient the edges, and enumerate the chords by $e_1,\ldots,e_b$.  Let
\[
 H=\Z\mathbf 1[0]\oplus L[1]
\]
with zero differential, identifying a basis $\ell_1,\ldots,\ell_b$ of $L$ with the fundamental cycles of the chords.  Define a cellular chain map
\begin{equation}\label{eq:pi}
 \pi:C_*(G;\Z)\longrightarrow H,
 \qquad
 \pi_0(v)=\mathbf1,
 \quad \pi_1(e_a)=\ell_a,
 \quad \pi_1(t)=0\ (t\in E(T)).
\end{equation}
It is a chain map because both endpoints of every edge map to $\mathbf1$.  It is an isomorphism on $H_0$.  On $H_1$, it records the chord coefficients; the fundamental cycles have chord-coefficient matrix equal to the identity.  Hence $\pi$ is an integral quasi-isomorphism.

For an intersection whose free-coordinate set is $F$, take the ordered tensor product $\pi_F=\bigotimes_{j\in F}\pi$.  Since the complexes are bounded and degreewise free, $\pi_F$ is a quasi-isomorphism.  These maps also commute strictly with every horizontal map in~\eqref{eq:cech}.  Indeed, removing one intersection requirement makes some formerly constrained coordinates free; the corresponding inclusion inserts their anchor vertices.  Formula~\eqref{eq:pi} sends every inserted vertex to the same degree-zero element $\mathbf1$.

We therefore obtain a map from $D$ to a finite double complex $E$ whose intersection term is $H^{\otimes F}$ and whose horizontal maps insert $\mathbf1$ in newly free coordinates.  It is a quasi-isomorphism in every vertical column.  The vertical filtrations of the bounded total complexes consequently give
\begin{equation}\label{eq:total-qiso}
 H_*(\Tot D)\cong H_*(\Tot E).
\end{equation}

The vertical differential of $E$ is zero.  Expand each tensor factor as $H=\Z\mathbf1\oplus L[1]$.  A summand is specified by the subset $J\subseteq[n]$ of coordinates carrying an $L$ factor.  It appears in an intersection precisely when no anchor requirement uses a coordinate of $J$.  Thus, for fixed $J$, its horizontal complex is the simplicial chain complex of the compatible injections into $[n]\setminus J$, tensored with $L^{\otimes |J|}$ and shifted by $|J|$.  After the order-preserving identification $[n]\setminus J\cong[n-|J|]$, this is the chain complex of $N_{r,n-|J|}$.  There are no such injections when $|J|>n-r$.  Lemma~\ref{lem:augmentation} and~\eqref{eq:total-qiso} now prove Theorem~\ref{thm:decomposition}.

\section{Reduction of the nerve to a tree}

\begin{lemma}\label{lem:nerve-free}
For every $m\geq r\geq1$, all groups $H_*(N_{r,m};\Z)$ are free abelian.
\end{lemma}

\begin{proof}
For $r=1$, every collection of injections is compatible, so $N_{1,m}$ is a simplex.  Suppose $r\geq2$.  Choose any finite tree $T_r$ containing $r$ distinguished vertices, denoted again by $K$.  Cover $\Sigma(T_r,K,m)$ by the subcomplexes $U_f$ indexed by $f\in\inj(K,[m])$ as above.  Its compatibility nerve is exactly $N_{r,m}$.

Every nonempty intersection is a product of trees and hence contractible.  Apply Lemma~\ref{lem:augmentation} to this cover.  In the resulting double complex, augmenting the cellular chains of each nonempty intersection to $\Z$ is a columnwise quasi-isomorphism, compatible with all inclusions.  Its target is the simplicial chain complex of $N_{r,m}$, concentrated in vertical degree zero.  The vertical filtration therefore yields
\[
 H_*(N_{r,m};\Z)\cong H_*(\Sigma(T_r,K,m);\Z).
\]
Kozlov's Theorem~2.5 states that the latter space is homotopy equivalent to a wedge of $(m-r)$-spheres; importantly, the anchor set in that theorem is an arbitrary subset of the tree's vertices~\cite[Definition~2.4 and Theorem~2.5]{kozlov2023}.  Its integral homology is thus free.
\end{proof}

For connected $G$ and nonempty $K$, Theorem~\ref{thm:decomposition} and Lemma~\ref{lem:nerve-free} prove freeness because $L$ is free and all sums are finite.  If $K=\varnothing$, then $\Sigma(G,\varnothing,n)=G^n$; the cellular homology of a graph is free, and the integral K\"unneth theorem gives the same conclusion for the product~\cite[Theorem~3B.6]{hatcher2002}.

\section{Disconnected graphs and the cograph question}

Write a nonempty finite graph as $G=\bigsqcup_{a\in A}G_a$, with each $G_a$ connected, and set $K_a=K\cap V(G_a)$.  The product $G^n$ is the finite disjoint union indexed by maps $\alpha:[n]\to A$.  Put $I_a=\alpha^{-1}(a)$.  The part of the anchored space in this component is empty unless $|I_a|\geq |K_a|$ for every $a$; otherwise, after reordering labelled coordinates, it is
\begin{equation}\label{eq:components}
 \prod_{a\in A}\Sigma(G_a,K_a,|I_a|).
\end{equation}
Every factor has free homology by the connected case, including the convention $\Sigma(G_a,\varnothing,0)=\{*\}$.  The integral K\"unneth theorem makes the homology of~\eqref{eq:components} free, and a finite disjoint union gives a finite direct sum.  If $G$ is empty, then $K=\varnothing$; the space is a point for $n=0$ and empty for $n>0$, again with free homology.  This completes the proof of Theorem~\ref{thm:main}.

\begin{corollary}
There is no cograph tuple $(G,K,n,i)$ with $n\geq|K|$ for which $\operatorname{Tor}H_i(\Sigma(G,K,n);\Z)$ is nonzero.
\end{corollary}

\section{Exact checks and scope}

The proof above is independent of computation.  Two accompanying Python implementations were used only to test cell conventions, signs, and rank consequences.  They construct cubical boundaries with integer coefficients, check $\partial^2=0$, and compare ranks over $\mathbb F_2$, $\mathbb F_3$, and $\mathbb F_5$.  The independent implementation exhausts all connected labelled simple graphs on at most four vertices: 167 graph/one-anchor cases at $n=3$ and 241 graph/two-anchor cases at $n=4$.  Every case agrees with~\eqref{eq:decomposition}.  Exact Smith computations additionally give
\[
 (\beta_0,\beta_1,\beta_2)=(1,9,27)
 \quad\text{for }(K_4,\{0\},3)
\]
and $(1,12,121)$ for $(K_4,\{0,1\},4)$, with all nonzero Smith factors equal to one.  These are finite consistency checks, not a proof of universal torsion-freeness.

For methodological comparison, Zeng's exact graph enumeration~\cite{zeng2026whiskered} likewise uses canonical graph representatives, exact integer arithmetic, and an independently checkable verifier that reconstructs matrices from graph encodings.  Here the two scripts reconstruct the anchored cubical chains independently; neither a bounded census nor agreement across several fields is promoted to an integral theorem.

The literature comparison for this note inspected the problem source and the cited tree and circle results through September~8, 2026.  It found no equivalent chain-level decomposition in those sources; that bounded observation is not evidence of exhaustive priority.

\section*{Computational and AI assistance}

OpenAI Codex assisted with proof exploration, implementation of the exact checks, verification planning, and LaTeX drafting.  Its output was not treated as mathematical evidence.  The chain argument was reconstructed in a separate audit, and the finite checks were rerun through two independently written implementations.  All reported computations use integer or finite-field arithmetic, with no floating-point inference.  The author reviewed the manuscript and is responsible for its claims.

\small
\bibliographystyle{plain}
\bibliography{references}

\end{document}
